\emph{(i)} $n$ lies in $\R^{1,3}$, which is orthogonal to $\nO$ and $\nInf$, so $\ip{n}{I} = 0$ and $Q(n) = 0$. For $x \in \R^{1,3}$, $\ip{X(x)}{n} = \ip{x}{n}$, so $N_{n} = \{x_{\perp} + s\,n\}$ is the background null plane with normal $n$; since $\hb^{-1}(n) = n/\sigma$ for the field $\hb = \sigma$ of Lemma \claimref{mass.de-sitter-scale-frozen-point}, the normal is physically null (Fact \claimref{synthesis.imported-gtg-dictionary}(i)). (Beyond one flat chart: $n$ is then a null vector of $I^{\perp}$, a point of the conformal boundary $\{\ip{Y}{I} = 0\}$ of the de Sitter field, and $N_{n}$ is its light cone.) There $Q(x) = Q(x_{\perp} + s\,n) = Q(x_{\perp}) + 2s\ip{x_{\perp}}{n} + s^{2}Q(n) = -|x_{\perp}|^{2}$, so $\sigma = 1 + \tfrac{\Lambda}{12}|x_{\perp}|^{2} > 0$, independent of $s$. The bivector $n \wedge m$ lies in $\Lambda^{2}\R^{1,3}$, which commutes with $\nO$ and $\nInf$ and hence with $I$; its rotor is a Lorentz boost of the background, and since $\sigma$ depends on $x$ only through $Q(x)$ the boost, with its induced rotation gauge, leaves $(\hb, \Omega)$ unchanged: it is a symmetry of the gravitational field, and its Killing field in the sense of Fact \claimref{synthesis.imported-gtg-dictionary}(vi) is normal to $N_{n}$ and null there. It sends $n \mapsto e^{-\theta}n$, $m \mapsto e^{\theta}m$ and fixes $x_{\perp}$, so it maps $N_{n}$ to itself, rescaling $s$ about $s = 0$: $N_{n}$ is a sheet of its Killing horizon with edge $\{s = 0\}$. The check verifies these identities exactly for $n = e_{0} + e_{1}$, $m = \tfrac12(e_{0} - e_{1})$, at $\Lambda = 3, \tfrac37, 12$ and three transverse points.

\emph{(ii)} The rotation-gauge field of $\hb = \sigma$ is $\Omega(a) = -(a \wedge \nabla\sigma)/\sigma$ (Lemma \claimref{mass.de-sitter-scale-frozen-point}), with $\nabla\sigma = -\tfrac{\Lambda}{6}x$. Along $x(\tau) = x_{\perp} + c\tau\,n$ the physical tangent is $v = \hb^{-1}(x') = c\,n/\sigma(x_{\perp})$, constant, and $\Omega(x')\cdot v = -\tfrac{c^{2}}{\sigma^{2}}(n \wedge \nabla\sigma)\cdot n = -\tfrac{c^{2}}{\sigma^{2}}\,n\,(n\cdot\nabla\sigma) = 0$ because $n \cdot \nabla\sigma = -\tfrac{\Lambda}{6}\ip{n}{x} = 0$ on $N_{n}$; so $dv/d\tau = -\Omega(x')\cdot v$ holds with $Q(v) = 0$, and by Fact \claimref{synthesis.imported-gtg-dictionary}(iii) these lines are affinely parametrised free-fall null worldlines for every $c$. The physical affine parameters are the affine functions of $s$, a choice on each generator; for $\lambda = s/\sigma(x_{\perp})$, $dx/d\lambda = \sigma\,n$ and $\hb^{-1}(dx/d\lambda) = n$ on every generator. $W$ is orthogonal to $I$, to $\R^{1,3}$ and so to $n$ and $m$, with $Q(W) = -\Lambda/3$; the bivectors $e_{2}e_{3}$, $e_{2}W$, $e_{3}W$ (for $n = e_{0} + e_{1}$) are built from vectors orthogonal to $I$, $n$ and $m$, so they commute with $I$ and with $n \wedge m$: acting on the conformal model they fix $I$, hence $\sigma$ up to the induced position gauge, and generate the rotations of the static patch of that boost, symmetries of the gravitational field fixing $n$. For $Y$ on the light cone of $n$, written as a multiple of $X(x_{\perp} + s\,n)$, the ratio $\ip{Y}{m}/\ip{Y}{I}$ does not depend on the multiple and equals $s/\sigma(x_{\perp}) = \lambda$, since $\ip{X(x_{\perp} + s\,n)}{m} = \ip{x_{\perp} + s\,n}{m} = s$; the rotations fix $m$ and $I$, so they preserve $\lambda$. They act transitively on the transverse points, so any rotation-invariant normalisation of the affine parameters is a constant multiple of $\lambda$. The check confirms the invariance for the rotors $1 + \tfrac13 B$ of the three bivectors $B$, at three transverse points and $s = 0, 1, \tfrac52$, with the $(e_{2}, W)$ images tabulated. The uniform affine translations are therefore $\lambda \mapsto \lambda + a$ on every generator.

\emph{(iii)} $n$ and $I$ are orthogonal, so $nI = n \wedge I = -In$ and $(nI)^{2} = -Q(n)Q(I) = 0$; hence $\exp(\tfrac{a}{2}nI) = 1 + \tfrac{a}{2}nI = T_{a}$, and $T_{a}\tilde T_{a} = 1$ (the null case of the exponential of a bivector with scalar square, \claimref{std-algebra.bivector-exponential}, which also gives $(n \wedge I)^{2} = \ip{n}{I}^{2} - Q(n)Q(I)$). Expanding $I$, $\tfrac12 nI = \tfrac12 n\nInf + \tfrac{\Lambda}{12}n\nO = P_{n} + \tfrac{\Lambda}{12}n\nO$. Acting by $v \mapsto T_{a}v\tilde T_{a}$, $T_{a}$ fixes $n$ and, for $Y$ with $\ip{Y}{n} = 0$, sends $Y \mapsto Y + a\ip{Y}{I}\,n$ (the terms of $\tfrac{a}{2}(nIY - YnI)$ and of second order reduce to this because $Y \perp n$ and $n$ is null); for $Y = X(x_{\perp} + s\,n)$, $\ip{Y}{I} = \sigma(x_{\perp})$, so $Y \mapsto X(x_{\perp} + s\,n) + a\,\sigma(x_{\perp})\,n = X(x_{\perp} + (s + a\sigma(x_{\perp}))n)$, the last because $Q$ is unchanged along the generator. So $\lambda$ advances by $a$ on every generator. The check confirms this exactly at $\Lambda = 3, \tfrac37$, $a = \tfrac25, -3$, three transverse points and two values of $s$, and that $P_{n}$ alone advances $\lambda$ by $1/\sigma(x_{\perp})$, which takes different values on different generators. $T_{a}$ moves $I$: $T_{a}I\tilde T_{a} = I + a\,Q(I)\,n \ne I$; so it is not a symmetry of the gravitational field: at $\Lambda = 3$, $a = \tfrac34$ it carries $X(e_{0})$, where $\sigma = \tfrac34$, to a vector $Y$ with $\ip{Y}{I} = 0$, a point where $\sigma = 0$ and $\hb$ is not invertible. It is a background conformal transformation that maps $N_{n}$ to itself and preserves $\sigma$ there, but not the frozen point.

\emph{(iv)} The cone $C_{\pi}$ is closed, convex and invariant under the adjoint action (Proposition \claimref{synthesis.positive-conformal-energy-gives-positive-null-translations}(ii)). From that result's orbit identity for the dilation of the chart of Lemma \claimref{mass.chart-form-spin-isomorphism}, $\mathrm{Ad}_{D_{\rho}}Z = \rho^{2}P_{e_{0}} + \rho^{-2}e_{0}\nO$, so $\rho^{2}\mathrm{Ad}_{D_{\rho}}Z = \rho^{4}P_{e_{0}} + e_{0}\nO \in C_{\pi}$ for every $\rho > 0$, and closedness gives $e_{0}\nO \in C_{\pi}$ as $\rho \to 0$. The boost $B_{\rho}$ of the stored result fixes $\nO$ and sends $e_{0} \mapsto \cosh\lambda\,e_{0} + \sinh\lambda\,e_{1}$ with $e^{\lambda} = \rho^{2}$, so $2\rho^{-2}\mathrm{Ad}_{B_{\rho}}(e_{0}\nO) = (e_{0} + e_{1})\nO + \rho^{-4}(e_{0} - e_{1})\nO$, and $(e_{0} + e_{1})\nO \in C_{\pi}$ in the limit; rotations fixing $e_{0}$, $\nO$ and $\nInf$ carry it to $(e_{0} + \hat n)\nO$ for every unit $\hat n$, and positive scaling gives $n\nO$ for every future null $n$. With $P_{n} \in C_{\pi}$ (the stored result (ii)) and convexity, $P_{n} + \tfrac{\Lambda}{12}n\nO \in C_{\pi}$ for $\Lambda \ge 0$. The check verifies the two orbit identities exactly, at $\rho = \tfrac12, \tfrac14, 3$ and $\rho = 2, 3$. For covariance, the static patches are permuted transitively by the identity component of the symmetry group of the gravitational field, the rotors of $\Cl(2,4)$ fixing $I$, which preserves $I$ and so carries $\tfrac12 nI$ to $\tfrac12 n'I$ for the image $n'$ of $n$, and $C_{\pi}$ is invariant; at $\Lambda = 3$ the check verifies that $R = 1 - e_{1}e_{W}$ has $R\tilde R = 2$, fixes $I$, and carries $n = e_{0} + e_{1}$ to $2v_{\Lambda}^{*}$, $v_{\Lambda}^{*} = \tfrac12(e_{0} - e_{W})$, the null vector of the static patch of the observer at $x = 0$ (Proposition \claimref{synthesis.cosmological-horizon-completed-pair}, whose static boost $e_{0}W$ is Proposition \claimref{synthesis.de-sitter-static-boost-frozen-point}). Beyond one flat chart: some of these rotors, and the static-patch rotations through $W$ in (ii), move points out of the chart, so the transitivity uses the compactified model. For $\Lambda < 0$ the coefficient of $n\nO$ is negative and the argument gives nothing, as it should: there is no cosmological horizon.

\emph{(v)} If the GNS vector $\Omega$ is invariant under the one-parameter group $\exp(t\cdot\tfrac12 nI)$, then by (iv) it is a ground state for that group, which by (iii) is the uniform affine translation of the generators of $N_{n}$, and by covariance for that of every cosmological horizon. A vector invariant under the translations of the chart is invariant under all of $\mathrm{Spin}^{+}(2,4)$ by the Howe-Moore theorem: in a continuous unitary representation of a connected simple Lie group with finite centre, a vector invariant under a closed noncompact subgroup is invariant under the group; $\mathrm{Spin}^{+}(2,4) \cong SU(2,2)$ has centre $\Z_{4}$ and the translations form such a subgroup (cited from memory; for the vacuum of a conformally invariant field it is the definition). The identification with case (b) of Premise \claimref{synthesis.n-plus-null-ground-state} needs the cosmological horizon's modes to be the restriction to $N_{n}$ of the field: for a conformally invariant field of conformal weight $w$ coupled to $(\hb, \Omega)$ with $\hb = \sigma$, the field is $\sigma^{w}$ times the field coupled to the vanishing gravitational field, and $\sigma$ is constant on each generator, so on each generator the two differ by a constant factor and the CAR normalisation in the parameter $\lambda$ fixes it; this is argued and not derived, and it holds only for conformally invariant fields (not the minimally coupled massless scalar). The positivity does not come from a conjugation: $\tfrac12 nI$ is not conjugate to a null translation, since $\mathrm{span}\{n, I\}$ is not totally null, and convexity of $C_{\pi}$ is what places it in the cone. \emph{(vi)} For a gravitational field of the form $\mathcal{R}(B) = cB$, as de Sitter's with $c = \Lambda/3$, the contraction is $\mathcal{R}(b) = \partial_{a}\cdot(c\,a \wedge b) = 3c\,b$, $\mathcal{R} = 12c$, $\Sch(a) = \tfrac{c}{2}a$, and the Weyl part of Fact \claimref{synthesis.imported-gtg-dictionary}(v) is $c\,a \wedge b - \tfrac{c}{2}a \wedge b - \tfrac{c}{2}a \wedge b = 0$; by Theorem \claimref{synthesis.gtg-tractor-gauge-field-reproduces-tractor-connection} the $\Cl(2,4)$ tractor gauge field, whose field strength on an Einstein field (Cotton term zero) is minus the Weyl part, then has zero field strength, and the background $\Cl(2,4)$ rotors used in (iii) and (iv) act on the field's null structure. For the Schwarzschild field, $\mathcal{R}(B) = -\tfrac{M}{r^{3}}S(B)$, $S = 1 - 3\Pi$, $\Pi$ the projection onto $\mathrm{span}\{\sigma_{r}, I\sigma_{r}\}$ (Proposition \claimref{singularity.curvature-gauge-field-exact-shape-decoupling}); in an orthonormal physical frame with $\sigma_{r} = e_{0}e_{1}$ the contraction $\partial_{a}\cdot S(a \wedge b)$ vanishes for every $b$ (computed exactly in the check), so $\Sch = 0$ and the Weyl part is $\mathcal{R}$ itself, nonzero at every $r > 0$ (spectrum $-\tfrac{M}{r^{3}}(-2, -2, 1, 1, 1, 1)$). Under a rotation gauge $\mathcal{R}(B) \mapsto U\mathcal{R}(\tilde UBU)\tilde U$ and under a displacement it transforms tensorially (Definition \claimref{obj.synthesis.gtg-flat-spacetime-gauge-fields}(e)), and the Weyl part with it; so no gauge transformation brings the Schwarzschild field to a pure dilation $\hb = \sigma$, whose Weyl part vanishes (for de Sitter by the computation above; for general $\sigma$ the Weyl part of the torsion-free field $\Omega(a) = -(a \wedge \nabla\sigma)/\sigma$ vanishes, standard, cited from memory). The argument of (iii) has no counterpart at the Schwarzschild horizon. Literature, from memory: a gravitational field with nonvanishing Weyl part has fewer than fifteen conformal symmetries, and the Schwarzschild field's are its four Killing symmetries, none of which acts on the horizon's generators as a uniform affine translation (the static one acts as a boost about the edge).
